% GENERATED FILE. Edit canonical lessons, then run npm run generate:books.
% canonical-source-hash: fnv1a64:16bb395f8ea6a87c
% canonical-lessons: JA-C449-yakedo, JA-C449-nenza, JA-C449-kossetsu, JA-C449-kizuato, JA-C449-nuguu, JA-R449-review-444-449

\chapter{Burns, Sprains and Scars}
\label{ch:ja-burns-sprains-and-scars}
\hlchaptermodality{449}

\begin{chapteropening}
\textbf{By the end of this chapter:} \emph{I can talk about burns, sprains and scars.}

Complete the last lesson of chapter 449: I can talk about burns, sprains and scars.
\end{chapteropening}

\section[\texorpdfstring{\ja{やけど} (yakedo)}{やけど (yakedo)}]{\ja{やけど} (yakedo) — a burn}

\label{lesson:JA-C449-yakedo}



\begin{quote}
Before the new one: say the Japanese for to go numb, then the Japanese for to bend.
\end{quote}

\subsection*{You'll want to know: \ja{やけど}}
\textbf{\ja{やけど}} — \emph{yakedo} — \textquotedblleft{}a burn\textquotedblright{}.

\begin{quote}
\textbf{\ja{なべで} \ja{てに} \ja{やけどを} \ja{しました。}} — \emph{nabe de te ni yakedo o shimashita} — \textquotedblleft{}I burned my hand on the pot.\textquotedblright{}
\end{quote}

\begin{cousinweb}
\textbf{\ja{やけど}} (\emph{yakedo}) is a native Japanese word (wago), inherited rather than borrowed from Chinese. JMdict marks it as usually written in kana; its usual kanji spelling writes the meaning, \textquotedblleft{}fire\textquotedblright{} and \textquotedblleft{}wound\textquotedblright{}, not the sound (JMdict, KANJIDIC2).
\end{cousinweb}

\begin{grammarlens}[title={What you've built}]
One more word for describing what you lived and felt.
\end{grammarlens}

\subsection*{Your turn}
\begin{itemize}
\raggedright
  \item \emph{Say it:} \emph{yakedo}
  \item \emph{Say it:} \emph{yakedo}, once more
  \item \emph{Say it:} \emph{nabe de te ni yakedo o shimashita} — I burned my hand on the pot.
  \item \emph{Recall:} say the Japanese for to go numb, then the Japanese for to bend, then say \emph{yakedo} again
\end{itemize}

\begin{tcolorbox}[breakable,colback=teal!4,colframe=teal!35!black,title={Before you move on}]
What does \ja{やけど} mean? (\textquotedblleft{}A burn\textquotedblright{}.) Say it once more.
\end{tcolorbox}
\section[\texorpdfstring{\ja{ねんざ} (nenza)}{ねんざ (nenza)}]{\ja{ねんざ} (nenza) — a sprain}

\label{lesson:JA-C449-nenza}



\begin{quote}
Before the new one: say the Japanese for to bend, then the Japanese for a burn.
\end{quote}

\subsection*{You'll want to know: \ja{ねんざ}}
\textbf{\ja{ねんざ}} — \emph{nenza} — \textquotedblleft{}a sprain\textquotedblright{}.

\begin{quote}
\textbf{\ja{しあいで} \ja{あしくびを} \ja{ねんざしました。}} — \emph{shiai de ashikubi o nenza shimashita} — \textquotedblleft{}I sprained my ankle in the match.\textquotedblright{}
\end{quote}

\begin{cousinweb}
\textbf{\ja{ねんざ}} (\emph{nenza}) is Sino-Japanese (kango): two Chinese characters read with their on'yomi, readings borrowed from Chinese. They are \emph{nen}, \textquotedblleft{}twist\textquotedblright{} (Middle Chinese \emph{nep}) and \emph{za}, \textquotedblleft{}sprain\textquotedblright{} (\emph{tswaH}) (KANJIDIC2; Middle Chinese in Baxter's transcription).
\end{cousinweb}

\begin{grammarlens}[title={What you've built}]
One more word for describing what you lived and felt.
\end{grammarlens}

\subsection*{Your turn}
\begin{itemize}
\raggedright
  \item \emph{Say it:} \emph{nenza}
  \item \emph{Say it:} \emph{nenza}, once more
  \item \emph{Say it:} \emph{shiai de ashikubi o nenza shimashita} — I sprained my ankle in the match.
  \item \emph{Recall:} say the Japanese for to bend, then the Japanese for a burn, then say \emph{nenza} again
\end{itemize}

\begin{tcolorbox}[breakable,colback=teal!4,colframe=teal!35!black,title={Before you move on}]
What does \ja{ねんざ} mean? (\textquotedblleft{}A sprain\textquotedblright{}.) Say it once more.
\end{tcolorbox}
\section[\texorpdfstring{\ja{こっせつ} (kossetsu)}{こっせつ (kossetsu)}]{\ja{こっせつ} (kossetsu) — a fracture; a broken bone}

\label{lesson:JA-C449-kossetsu}



\begin{quote}
Before the new one: say the Japanese for a burn, then the Japanese for a sprain.
\end{quote}

\subsection*{You'll want to know: \ja{こっせつ}}
\textbf{\ja{こっせつ}} — \emph{kossetsu} — \textquotedblleft{}a fracture; a broken bone\textquotedblright{}.

\begin{quote}
\textbf{\ja{ころんで、うでを} \ja{こっせつしました。}} — \emph{koronde, ude o kossetsu shimashita} — \textquotedblleft{}I fell over and broke my arm.\textquotedblright{}
\end{quote}

\begin{cousinweb}
\textbf{\ja{こっせつ}} (\emph{kossetsu}) is Sino-Japanese (kango): two Chinese characters read with their on'yomi, readings borrowed from Chinese. They are \emph{ko}, \textquotedblleft{}bone\textquotedblright{} (Middle Chinese \emph{kwot}) and \emph{setsu}, \textquotedblleft{}break\textquotedblright{} (\emph{tsyet}) (KANJIDIC2; Middle Chinese in Baxter's transcription).
\end{cousinweb}

\begin{grammarlens}[title={What you've built}]
One more word for describing what you lived and felt.
\end{grammarlens}

\subsection*{Your turn}
\begin{itemize}
\raggedright
  \item \emph{Say it:} \emph{kossetsu}
  \item \emph{Say it:} \emph{kossetsu}, once more
  \item \emph{Say it:} \emph{koronde, ude o kossetsu shimashita} — I fell over and broke my arm.
  \item \emph{Recall:} say the Japanese for a burn, then the Japanese for a sprain, then say \emph{kossetsu} again
\end{itemize}

\begin{tcolorbox}[breakable,colback=teal!4,colframe=teal!35!black,title={Before you move on}]
What does \ja{こっせつ} mean? (\textquotedblleft{}A fracture; a broken bone\textquotedblright{}.) Say it once more.
\end{tcolorbox}
\section[\texorpdfstring{\ja{きずあと} (kizuato)}{きずあと (kizuato)}]{\ja{きずあと} (kizuato) — a scar}

\label{lesson:JA-C449-kizuato}



\begin{quote}
Before the new one: say the Japanese for a sprain, then the Japanese for a fracture.
\end{quote}

\subsection*{You'll want to know: \ja{きずあと}}
\textbf{\ja{きずあと}} — \emph{kizuato} — \textquotedblleft{}a scar\textquotedblright{}.

\begin{quote}
\textbf{\ja{ひざに} \ja{ちいさい} \ja{きずあとが} \ja{あります。}} — \emph{hiza ni chiisai kizuato ga arimasu} — \textquotedblleft{}I have a small scar on my knee.\textquotedblright{}
\end{quote}

\begin{cousinweb}
\textbf{\ja{きずあと}} (\emph{kizuato}) is a native Japanese compound (wago), inherited rather than borrowed from Chinese: \emph{kizu}, \textquotedblleft{}wound\textquotedblright{} and \emph{ato}, \textquotedblleft{}mark\textquotedblright{} (JMdict, KANJIDIC2).
\end{cousinweb}

\begin{grammarlens}[title={What you've built}]
One more word for describing what you lived and felt.
\end{grammarlens}

\subsection*{Your turn}
\begin{itemize}
\raggedright
  \item \emph{Say it:} \emph{kizuato}
  \item \emph{Say it:} \emph{kizuato}, once more
  \item \emph{Say it:} \emph{hiza ni chiisai kizuato ga arimasu} — I have a small scar on my knee.
  \item \emph{Recall:} say the Japanese for a sprain, then the Japanese for a fracture, then say \emph{kizuato} again
\end{itemize}

\begin{tcolorbox}[breakable,colback=teal!4,colframe=teal!35!black,title={Before you move on}]
What does \ja{きずあと} mean? (\textquotedblleft{}A scar\textquotedblright{}.) Say it once more.
\end{tcolorbox}
\section[\texorpdfstring{\ja{ぬぐう} (nuguu)}{ぬぐう (nuguu)}]{\ja{ぬぐう} (nuguu) — to wipe away}

\label{lesson:JA-C449-nuguu}



\begin{quote}
Before the new one: say the Japanese for a fracture, then the Japanese for a scar.
\end{quote}

\subsection*{You'll want to know: \ja{ぬぐう}}
\textbf{\ja{ぬぐう}} — \emph{nuguu} — \textquotedblleft{}to wipe away\textquotedblright{}.

\begin{quote}
\textbf{\ja{かれは} \ja{あせを} \ja{ぬぐいました。}} — \emph{kare wa ase o nuguimashita} — \textquotedblleft{}He wiped away the sweat.\textquotedblright{}
\end{quote}

Polite \textbf{\ja{ぬぐいます}}, \emph{nuguimasu}; a godan verb. Past (polite): \textbf{\ja{せんしゅが} \ja{なみだを} \ja{ぬぐいました。}}, \emph{senshu ga namida o nuguimashita}.

\begin{cousinweb}
\textbf{\ja{ぬぐう}} (\emph{nuguu}) is a native Japanese word (wago), inherited rather than borrowed from Chinese; its kanji spelling was attached to it later (JMdict).
\end{cousinweb}

\begin{grammarlens}[title={What you've built}]
One more word for describing what you lived and felt.
\end{grammarlens}

\subsection*{Your turn}
\begin{itemize}
\raggedright
  \item \emph{Say it:} \emph{nuguu}
  \item \emph{Say it:} \emph{nuguu}, once more
  \item \emph{Say it:} \emph{kare wa ase o nuguimashita} — He wiped away the sweat.
  \item \emph{Recall:} say the Japanese for a fracture, then the Japanese for a scar, then say \emph{nuguu} again
\end{itemize}

\begin{tcolorbox}[breakable,colback=teal!4,colframe=teal!35!black,title={Before you move on}]
What does \ja{ぬぐう} mean? (\textquotedblleft{}To wipe away\textquotedblright{}.) Say it once more.
\end{tcolorbox}
\section[(dialogue)]{Looking back: from Breathing and Pulse to Burns, Sprains and Scars}

\label{lesson:JA-R449-review-444-449}



\begin{quote}
Say the Japanese for to breathe and a pulse.
\end{quote}

\subsection*{Your turn}
Say these aloud:

\begin{itemize}
\raggedright
  \item to breathe, a pulse, blood pressure, one's weight, snoring
  \item a headache, a stomach ache, throbbing, nausea, a bruise
  \item a sneeze, a runny nose, a chill, gargling, physical condition
  \item dizziness, unsteadily, sluggish, worn out, lack of sleep
  \item to get stiff, to massage, to twist, to go numb, to bend
  \item a burn, a sprain, a fracture, a scar, to wipe away
\end{itemize}

\begin{tcolorbox}[breakable,colback=teal!4,colframe=teal!35!black,title={Before you move on}]
To breathe? (\textbf{\ja{こきゅうする}}, \emph{kokyū suru}.) A pulse? (\textbf{\ja{みゃく}}, \emph{myaku}.) Blood pressure? (\textbf{\ja{けつあつ}}, \emph{ketsuatsu}.) One's weight? (\textbf{\ja{たいじゅう}}, \emph{taijū}.) Snoring? (\textbf{\ja{いびき}}, \emph{ibiki}.) A headache? (\textbf{\ja{ずつう}}, \emph{zutsū}.) A stomach ache? (\textbf{\ja{ふくつう}}, \emph{fukutsū}.) Throbbing? (\textbf{\ja{ずきずき}}, \emph{zukizuki}.) Nausea? (\textbf{\ja{はきけ}}, \emph{hakike}.) A bruise? (\textbf{\ja{あざ}}, \emph{aza}.) A sneeze? (\textbf{\ja{くしゃみ}}, \emph{kushami}.) A runny nose? (\textbf{\ja{はなみず}}, \emph{hanamizu}.) A chill? (\textbf{\ja{さむけ}}, \emph{samuke}.) Gargling? (\textbf{\ja{うがい}}, \emph{ugai}.) Physical condition? (\textbf{\ja{たいちょう}}, \emph{taichō}.) Dizziness? (\textbf{\ja{めまい}}, \emph{memai}.) Unsteadily? (\textbf{\ja{ふらふら}}, \emph{furafura}.) Sluggish? (\textbf{\ja{だるい}}, \emph{darui}.) Worn out? (\textbf{\ja{ぐったり}}, \emph{guttari}.) Lack of sleep? (\textbf{\ja{ねぶそく}}, \emph{nebusoku}.) To get stiff? (\textbf{\ja{こる}}, \emph{koru}.) To massage? (\textbf{\ja{もむ}}, \emph{momu}.) To twist? (\textbf{\ja{ひねる}}, \emph{hineru}.) To go numb? (\textbf{\ja{しびれる}}, \emph{shibireru}.) To bend? (\textbf{\ja{まげる}}, \emph{mageru}.) A burn? (\textbf{\ja{やけど}}, \emph{yakedo}.) A sprain? (\textbf{\ja{ねんざ}}, \emph{nenza}.) A fracture? (\textbf{\ja{こっせつ}}, \emph{kossetsu}.) A scar? (\textbf{\ja{きずあと}}, \emph{kizuato}.) To wipe away? (\textbf{\ja{ぬぐう}}, \emph{nuguu}.)
\end{tcolorbox}
